\section{خلاصة الفصل }
شكل هذا الفصل الركيزة الأساسية لتحويل النماذج الرياضية النظرية (التي تم اشتقاقها في الفصول السابقة) إلى سلوك حركي منضبط وموجه للذراع الروبوتية في فضاء العمل. تم في هذا السياق تصميم واختبار بنية تحكم هرمية متكاملة تتألف من مستويين رئيسيين (عالي ومنخفض المستوى)، وذلك لضمان تتبع دقيق للمسارات الكارتزية والمفصلية على حد سواء مع الحفاظ على استقرار واستجابة النظام في ظروف التشغيل المختلفة. 

ويمكن تلخيص أبرز الاستنتاجات العلمية والنتائج التقنية التي تم التوصل إليها من خلال دراسة وتحليل خوارزميات التحكم المصممة في هذا الفصل وفق النقاط الموسعة التالية:

\begin{enumerate}
    \item \textbf{التهيئة الرياضية وصياغة فضاء الحالة غير الخطي:}
    أثبتت الدراسة أن عزل مصفوفة القصور الذاتي وتسارعات المفاصل $\ddot{q}$ من معادلة الحركة الديناميكية الأساسية، ودمجها مع ديناميكية المشغلات الكهربائية (المحركات) والمخفضات، هو خطوة لا غنى عنها للانتقال من النموذج الميكانيكي البحت إلى نموذج نظام ميكاترونيكي متكامل. وقد أتاحت صياغة النظام على هيئة فضاء حالة غير خطي موحد بالصيغة $\dot{x}=f(x,u)$ تمثيلاً رياضياً دقيقاً لأبعاد النظام (حيث تم أخذ زوايا المفاصل الخمسة وسرعاتها الزاوية كمتجهات حالة، وإشارات الجهد كمداخل تحكم)، مما وفر منصة متينة لبناء خوارزميات التحكم المتقدمة وتحليل استقرارها.

    \item \textbf{التحليل المقارن لأداء متحكمات المستوى المنخفض (LLC):}
    تم إجراء محاكاة لتقييم أداء ثلاث استراتيجيات تحكم مفصلية مختلفة تحت تأثير ظروف تشغيلية قياسية وعند إدخال ضجيج القياس (Measurement-Noise) لاختبار مدى متانة الخوارزميات، وأسفرت المقارنة عن الاستنتاجات التالية:
    \begin{itemize}
        \item \textbf{متحكم PID مع تعويض الجاذبية:} على الرغم من بساطة تركيبه وقدرة حد تعويض الجاذبية على إلغاء الأحمال الاستاتيكية الناتجة عن أوزان الوصلات، إلا أن الأداء الخطي للمتحكم عانى من قصور واضح عند تتبع المراجع الديناميكية السريعة. وعلاوة على ذلك، أثبتت النتائج أنه كان \textbf{الأكثر تأثراً وحساسية لضجيج القياس}، حيث تسبب حد الاشتقاق ($D$) في تضخيم الضجيج عالي التردد، مما انعكس سلباً على استقرار السرعات وخاصة عند المفاصل ذات الكسب المرتفع.
        \item \textbf{التحكم التكيفي المعتمد على العزم المحسوب (ACTC):} أظهرت هذه الخوارزمية تفوقاً ملحوظاً واستقراراً ديناميكياً عالياً؛ ويعود ذلك إلى استغلالها المباشر لمصفوفات النموذج الديناميكي ($M$ و $G$) لإلغاء اللاخطية في النظام، بالإضافة إلى قدرة آلية التكيف \LR{(Adaptation Mechanism)} على تحديث المعاملات ذاتياً أثناء التشغيل. وقد منحت هذه الخاصية المتحكم حصانة ممتازة ضد ضجيج القياس وقدرة عالية على تقليل خطأ التتبع في الحالة المستقرة.
        \item \textbf{التحكم بالنمط الانزلاقي (SMC):} أكدت النتائج المتانة العالية \LR{(High Robustness)} التي تتمتع بها هذه الخوارزمية في مواجهة الاضطرابات، إلا أن تطبيقها العملي يتطلب دقة هائلة في ضبط معاملات حد الانزلاق واختيار سمك "طبقة التشبع" \LR{(Boundary Layer)} بدقة، لكونها تمثل عامل الموازنة الحرج بين تقليل خطأ التتبع وتخفيف ظاهرة التشذر أو الاهتزاز عالي التردد \LR{(Chattering)} الضارة بالمشغلات الميكانيكية.
    \end{itemize}

    \item \textbf{تقييم تكامل حلقة التحكم عالي المستوى (HLC):}
    أظهرت نتائج محاكاة طبقة التحكم الكارتزي نجاحاً ملموساً في دمج الحلول الحركية الأمامية والعكسية مع حلقة التحكم المفصلي. وبيّنت الدراسة أن سلسلة التدفق المقترحة:
    \[
    X_{ref} \rightarrow \text{IK} \rightarrow q_{ref} \rightarrow \text{LLC} \rightarrow q \rightarrow \text{FK} \rightarrow X
    \]
    تتمتع باستجابة ديناميكية متوافقة، حيث تقاربت الإحداثيات الحقيقية لنهاية الذراع (End-Effector) مع القيم الكارتزية المرجعية المطلوبة بكفاءة عالية. وعلى الرغم من ظهور مرحلة انتقالية أولية \LR{(Transient Phase)} ناتجة عن القصور الذاتي للذراع وحركته الانتقالية من وضعه الابتدائي وتأثير مولد المراجع، إلا أن النظام حقق الاستقرار المطلوب \LR{(Steady-State Stability)} دون وجود اهتزازات مجددة.
\end{enumerate}

بناءً على هذه الاستنتاجات الشاملة، فإن نجاح الاستراتيجية الهرمية المصممة واستقرارها في بيئة المحاكاة الرقمية يثبت الجاهزية الرياضية والهندسية الكاملة للنظام.  